

\section{\texorpdfstring{Expected impact
}{Expected impact }}\label{expected-impact}

\subsection{Scale of PLAN4ARI contribution to the
outcomes}\label{scale-of-plan4ari-contribution-to-the-outcomes}

\textbf{Contribution to AI in Manufacturing.} The global AI in
manufacturing market size is projected to reach USD 9.89 Billion by 2027
(CAGR of 24.2\%).\footnote{\url{https://www.fortunebusinessinsights.com/artificial-intelligence-ai-in-manufacturing-market-102824}}
The WEF, instead, claims that the global AI in the manufacturing market
is poised to grow to \$20.8 billion by 2028.\footnote{\url{https://www.weforum.org/agenda/2024/01/how-we-can-unleash-the-power-of-ai-in-manufacturing/}}
McKinsey claims that an organization\textquotesingle s average number of
AI capabilities has doubled in the past five years. Still, many
organizations need help to maintain this momentum as they scale their AI
efforts\footnote{\url{https://www.mckinsey.com/capabilities/mckinsey-digital/our-insights/tech-forward/scaling-ai-for-success-four-technical-enablers-for-sustained-impact}}
because technical and cultural barriers at all hierarchical levels
hamper the introduction of AI-powered mobile robots. GMI\footnote{\url{https://www.gminsights.com/industry-analysis/artificial-intelligence-ai-in-manufacturing-market}}
underlines that the risk is that the Internet Giants will manage the
market, offering a powerful business development platform. These risks
are motivated by the fact that the machine learning techniques leading
AI in manufacturing need to generate and maintain enormous datasets,
which are often complex. Considering the EU\textquotesingle s digital
strategy\footnote{\url{https://digital-strategy.ec.europa.eu/en/news/investing-ai-manufacturing}},
the impact of AI actions relies on seizing AI opportunities to enhance
Europe\textquotesingle s mid and long-term competitiveness. In such a
scenario, \brand{Plan4ARI} implements a platform that may boost SMEs
that are 70\% of the expected market\footnote{\url{https://www.weforum.org/agenda/2024/01/how-we-can-unleash-the-power-of-ai-in-manufacturing/}}.
\brand{Plan4ARI} roadmap consists of the following\textbf{:}

\begin{enumerate}
\def\labelenumi{\arabic{enumi}.}
\item
  demonstrating that task planning can improve productivity
  \textbf{without a large amount of data}. This is achieved by the
  \textbf{task and motion planner that autonomously deploys the tailored
  solver for each application}.
\item
  allowing each company to exploit its resources without relying on
  external services. This is achieved by adequately designing the KB.
\item
  defining usage and design protocols that will pave the path to scaling
  AI-powered robots\footnote{Arinez, J. F., et al. "Artificial
    Intelligence in Advanced Manufacturing: Current Status and Future
    Outlook." ASME. J. Manuf. Sci. Eng. November 2020.}. The
  investigation of trustworthiness achieves this.
\end{enumerate}

The PLAN4ARI concept copes with at least 20\% of the tasks that can be
automated, so a market of \$5 billion by 2028. PLAN4ARI will exploit
available architectures to contribute to the EU AI ecosystem
proactively.

\textbf{Contribution to Mobile Manipulation in Manufacturing.} The high
flexibility of mobile manipulators is a boost for companies with a
``high mix low volume'' (HMLV). Remarkably, such industries display that
at least 85\% of production tasks may be automatized\footnote{M. Zinser,
  J. Rose, H. Sirkin, Industries and economies leading the robotics
  revolution, BCG Global (2015).}. The WEF Future of Jobs 2023 report
claims that the impact of such a technology on jobs will be positive
over the next five years\footnote{World Economic Forum, The future of
  jobs report 2023 (2023).}. Integrating mobile manipulators bears
complexity, high acquisition costs, and unsolved safety and security
challenges, and it is a challenge for SMEs\footnote{Jakob Gros, et al.
  Unlocking the Benefits of Mobile Manipulators for Small and
  Medium-Sized Enterprises, Procedia CIRP (120), 2023,}. Omron claims
that mobile manipulation may improve capacity gain by up to 8\%. Mobile
manipulators may be an essential value-added only if SMEs can use them
efficiently \textsuperscript{21}. The SOTA highlights that mobile
manipulators miss simple control frameworks more than advanced
functionalities. \brand{Plan4ARI} roadmap consists of the following:

\begin{enumerate}
\def\labelenumi{\arabic{enumi}.}
\setcounter{enumi}{3}
\item
  Having selected end-users that are paradigmatic examples of HMLV
  companies. Fa
\item
  adopting all the technologies that make the usage of mobile
  manipulators transparent.
\end{enumerate}

The consortium estimates that the \brand{Plan4ARI} approach will allow
coping with at least 25\% of the intralogistics market.

\textbf{Contribution to collaborative robotics in manufacturing}. The
number of collaborative robots installed has a share of 7.5\% in 2021.
EU robot installations grew by 24\% to 84k units in 2021, a new peak for
the region. \textasciitilde39k out of more than 517k robots installed
were cobots \footnote{\url{https://ifr.org/downloads/press2018/2022_WR_extended_version.pdf}}.
Although this market is proliferating, it is in the startup phase.
McKinsey estimates that about 30\% of the market is hampered by cultural
barriers that limit adoption\footnote{\url{https://www.mckinsey.com/capabilities/mckinsey-digital/our-insights/tech-forward/scaling-ai-for-success-four-technical-enablers-for-sustained-impact}}.
In such a scenario, the \textbf{PLAN4ARI} roadmap consists of the
following:

\begin{enumerate}
\def\labelenumi{\arabic{enumi}.}
\setcounter{enumi}{5}
\item
  developing tools that are easy and lightweight and fit
  SMEs\textquotesingle{} needs to guide the design of trustworthy cells
  accepted and used by human operators. Workers will evaluate the tools
  as detailed in Activity ID6.
\end{enumerate}

Consortium estimation is that this approach will boost the adoption of
collaborative robotics in at least 10\% of the companies that are
reluctant because of cultural barriers.

\subsection{Significance}\label{significance}

The \brand{Plan4ARI\textquotesingle s} demonstration that it is
possible to integrate AI into standard industrial toolchains will boost
the distribution of robotic technologies in industries where the added
value of final products, process timings, and working conditions may
vary significantly.

On the one hand, \brand{Plan4ARI} will

\begin{itemize}
\item
  \textbf{Reduce the errors.} According to industrial practice, up to
  5\% of tasks may have errors or defects in current manual solutions.
  Lean manufacturing improves the flow of operations, although it cannot
  nullify them. AI-powered robots integrated with holistic control
  enable the implementation of continuous quality improvement methods.
\item
  \textbf{Reduce the costs.} Optimizing the flow of operations through
  the plant and exploiting the human-robot hybrid approach makes it
  feasible to increase the number of orders the same number of employees
  can prepare.
\item
  \textbf{Reduce the space.} Manufacturers demand small-footprint
  solutions, as space-saving impacts the final cost.
\end{itemize}

On the other hand, \brand{Plan4ARI} will improve productivity in terms
of three KPIs. Precisely, it will improve:

\begin{itemize}
\item
  \textbf{The number of orders by the system.} The shopfloor
  orchestrator will allow a continuous reconfiguration, reducing the
  lead time from the order to the production start.
\item
  \textbf{The throughput.} The shopfloor orchestrator will optimize the
  shopfloor for productivity maximization.
\item
  \textbf{The ergonomics (physical and cognitive).} We will demonstrate
  that optimizing human trustworthiness corresponds to an ergonomics
  improvement.
\end{itemize}

\brand{Plan4ARI} will achieve these goals by implementing three
actions:

\begin{enumerate}
\def\labelenumi{\arabic{enumi}.}
\setcounter{enumi}{6}
\item
  implementing a platform-independent framework with little adaptation
  to specific hardware. This will allow the exploitation of existing HW
  and the selection of the HW setup based only on the needed
  performances.
\item
  Combining one or more collaborative modes safely depending on the
  specific end-user layout, maintaining the risk estimation and risk
  reduction measures concentrated at the end-user system integration.
\item
  To develop an easy-to-use setup with intuitive interfaces granting the
  usability requirements of the simplest expertise and reducing barriers
  to wide diffusion and major reconfiguration and control overhaul
\end{enumerate}

\subsection{Socio-economic Impact}\label{socio-economic-impact}

The goal is to achieve a 15\% increase in the OECD Job Quality Index
through work environment and safety improvement.

Job quality is a multi-dimensional concept that affects
workers\textquotesingle{} well-being. The OECD framework for measuring
and assessing job quality considers three objective and measurable
dimensions of job quality that are both important for worker well-being
and relevant for policy, such as the quality of the working environment.

Together, they provide a comprehensive assessment of job quality.
\textbf{Plan 4ARI aims to increase the job quality index by at least
15\% through safety improvement and a better} working environment.

One of the main impacts on worker safety will be the development of a
robust and precise system that can work at higher flexibility while
maintaining a high safety level. With robots learning about the physical
aspects of workers, e.g., their height, speed of task completion, motion
model, etc., they can provide better and personalized safety functions
for each worker. \textbf{This will also increase the level of trust that
the workers will have while collaborating with robots.} Furthermore,
increasing the worker\textquotesingle s cooperation experience and the
system\textquotesingle s productivity will create consistent work
environments that allow workers to be more productive and less stressed.
Through the \brand{Plan4ARI} solution, ergonomically appropriate
workplaces for human operators will be ensured. At the same time,
information and instructions concerning the process and production will
be provided to them highly intuitively. \brand{Plan4ARI} consortium
will study in detail the EU directive 89/391/EEC on measures to improve
safety and health at work. It will contribute to designing and deploying
working environments that will fulfill the existing expected
specifications.

\brand{Plan4ARI} will propose an innovative and flexible system that
involves all stakeholders in designing and organizing workspaces.
Indeed, to collaborate with robots, workers will be trained and acquire
high skills that will perform well in all dimensions. Returning to these
skills will result in higher employment and better jobs with lower
strain. In this context, \brand{Plan4ARI} will increase Job Quality by
at least 15\% by fostering workers\textquotesingle{} security and
creating comfortable working environments. The OECD database and studies
could be pertinent for estimating this index. Therefore,
\brand{Plan4ARI} will use the four potential OECD indicators of job
quality: time pressure, work autonomy, physical health risk factors, and
social support at work. To calculate the Job Quality Index through the
working environment, \brand{Plan4ARI} will define tailored surveys
based essentially on these four principal factors and propose questions
related to human-robot collaboration. Increasing job quality will reduce
costs resulting from workplace accidents and work-related health
problems, improve security, and enhance worker satisfaction. Moreover,
fewer accidents and work-related health problems relieve pressure on
public and private social protection, insurance, and pension systems.

\subsection{Requirements and potential barriers to achieving the
outcomes and impacts. Mitigation
measures}\label{requirements-and-potential-barriers-to-achieving-the-outcomes-and-impacts.-mitigation-measures}

\subsubsection{Technical Limitations}\label{technical-limitations}

Factor \#1 - Market Challenger. \brand{Plan4ARI\textquotesingle s}
holistic platform is unique, while new market incomers may challenge the
different AI-embedded controllers. If the challenger is a scientific
methodology not covered by IPR, the holistic platform is designed to
include any further solver. Conversely,
\brand{Plan4ARI\textquotesingle s} open architecture allows bridges to
proprietary tools without losing their potential. The route to
exploitation will leverage the Free and Open-Source Software (FOSS) - or
simply Open-Source Software (OSS) - tools, and the market outcome will
be a set of services.

Factor \#2 -- Performance. Many levels in architecture could lower the
overall performance. The fine-tuning of the system will allow for
identifying a sub-optimal solution in a limited time. As for Factor \#1,
the platform\textquotesingle s modularity will allow the rescaling and
reconnection of different solvers to improve the performance.

Factor \#3 -- Scalability. The devised architecture allows
straightforward scalability. \brand{Plan4ARI} will reduce the number of
dependencies to the minimum to allow users to select the necessary
modules. The dependencies will be done using an abstract interface to
allow the modules to be used without the explicit need for the whole
framework.

Factor \#4 - Computational burden. This technical limitation could arise
at many levels. Each module is designed so that it will be possible to
relax some optimality to grant at least feasibility.

\subsubsection{Economic Limitations}\label{economic-limitations}

Factor \#5 - Financing of follow-up steps. The business cases will
identify the financial needs and sources. For the results that adopt
open-source distribution, it will be possible to benefit from the
improvements and tests granted by the community.

\paragraph{Policy and Adoption
Limitations}\label{policy-and-adoption-limitations}

Factor \#6 - Legal framework. The legal framework will change during the
project time limit (new ISO10218). The Consortium will continue to
update the plan according to such innovations. The AI Act claims
limitations in tracking emotions in the workplace. Despite tracking the
humans inside the scene not being one-to-one with emotions, the risk
that such technologies could raise issues in future business deployment
is significant. \brand{Plan4ARI\textquotesingle s} roadmap to overcome
these limitations follows the GDPR approach: (i) the complete
anonymization of the biometrics tracked data; (ii) methodologies that do
not store biometrics data but that use this information online; (iii)
full ethics checkup of the sustainability of such methodologies.

Factor \#7 - Users\textquotesingle{} Acceptance of AI-powered mobile
robots. The tools and metrics will also allow us to understand how to
propose such technologies to standard end-users to avoid hype or
delusional rejection.